\section{Transfer Protocol for Other Parks}
The value of the framework is not restricted to Etosha. Its core logic is portable because it is built around operating roles rather than location-specific formulas. What transfers is the protected-area operating system; what changes are the data layers that populate it.

\subsection{Invariant System Components}
The following components transfer directly across reserves:
\begin{itemize}
  \item spatial control zoning,
  \item ecological consequence scoring,
  \item access- and pressure-based threat priors,
  \item constrained multi-asset allocation,
  \item closed-loop simulation and stress testing.
\end{itemize}

\subsection{Recalibration Protocol}
Transfer should follow a four-step protocol:
\begin{enumerate}
  \item rebuild the park digital twin with local geography, species, and mobility data;
  \item re-estimate threat priors using local access logic and intelligence patterns;
  \item regenerate patrol libraries using the local road, river, or trail network;
  \item re-identify the staffing floor and safe operating envelope under local response conditions.
\end{enumerate}
This protocol preserves the system architecture while correctly localizing its parameters.

\subsection{Illustrative Transfers}
For Yellowstone, winter closures, visitor pressure, and species-specific conflict patterns alter \(A_i\), \(D_i\), and seasonal multipliers \citep{yellowstoneOfficial}. For Kakadu, wet--dry flooding, river access, and habitat fragmentation reshape the mobility layer and shift the timing of exposure \citep{kakaduOfficial}. In both cases, the same controller logic remains valid, but the dominant bottlenecks move because the mobility system changes.

\begin{figure}[H]
  \centering
  \includegraphics[width=\linewidth]{fig11_transferability_matrix.png}
  \caption{Transfer protocol separating components that are structurally portable from those that require local rebuilding.}
  \label{fig:transfer}
\end{figure}
\FloatBarrier

\begin{table}[H]
  \centering
  \caption{Operating-system components: universal core versus park-specific recalibration.}
  \label{tab:transfer}
  \begin{tabular}{p{0.28\linewidth}p{0.24\linewidth}p{0.38\linewidth}}
    \toprule
    Component & Universal? & Required recalibration \\
    \midrule
    Zoning and risk architecture & Yes & Cell size and zone boundaries \\
    Composite priority logic & Yes & Weight values and seasonal modifiers \\
    Route-based controller & Yes & Local mobility network and station geometry \\
    Detection/interception simulation & Yes & Asset rates and response-delay profile \\
    Wildlife sensitivity layer & No & Species distribution and seasonal aggregation \\
    Intrusion prior & No & Local illegal-activity patterns and intelligence structure \\
    \bottomrule
  \end{tabular}
\end{table}

\subsection{Expected Strategic Shifts}
Under Yellowstone-like conditions, we expect stronger seasonal oscillation in \(\Resp\) because mobility is more volatile and visitor flows alter access patterns. Under Kakadu-like wet--dry transitions, the controller would shift toward season-adaptive route families, higher uncertainty buffers, and stronger emphasis on mobility rebuilding each season. In both settings, the underlying management lesson remains the same: staffing and hotspot revisit reliability dominate long-run protection quality.
